\documentclass[10pt,twocolumn,a4paper]{article}
\usepackage[T1]{fontenc}
\usepackage{lmodern}
\usepackage{amsmath}
\usepackage{graphicx}
\usepackage{pgf}
\usepackage{xcolor}
\usepackage{microtype}

\setlength{\columnsep}{0.23in}
\setlength{\parindent}{0pt}
\setlength{\parskip}{0.34em}
\setcounter{topnumber}{4}
\setcounter{totalnumber}{6}
\renewcommand{\topfraction}{0.92}
\renewcommand{\textfraction}{0.06}
\renewcommand{\floatpagefraction}{0.72}

\title{\textbf{Publication-ready figures}\\
\large LaTeX fonts, physical size, and accessibility}
\author{A reproducible \texttt{contrastcolors} demonstration}
\date{}

\begin{document}
\maketitle

\begin{abstract}
This two-column article demonstrates a document-aware Matplotlib workflow.
Figures are generated from dimensions and font sizes measured from this
LaTeX document. One graph is included using PGF, so its text is typeset
with the same LaTeX font configuration as the surrounding paragraphs.
The remaining examples illustrate a vector PDF, rasterized scatter artists,
and six simulated accessibility conditions.
\end{abstract}

\section{What does font matching mean?}
Scientific plots should be designed at their \emph{final physical width}.
The preamble probe measures \verb|\columnwidth|, \verb|\textwidth| and
normal body font size. In these plots, axis labels use the body point size;
tick and legend labels follow a consistent hierarchy.

A PDF produced independently by Matplotlib can have the correct
font \emph{size} without matching its glyphs.
PGF instead instructs LaTeX to typeset text when the figure is included.

\begin{figure}[tb]
  \centering
  \input{generated/lines.pgf}
  \caption{\textbf{Native PGF.} Text and math labels use the font
  configuration of this document. Markers and line patterns identify the
  plotted series in addition to their colors.}
  \label{fig:pgf}
\end{figure}

\begin{figure}[tb]
  \centering
  \includegraphics{generated/lines.pdf}
  \caption{\textbf{Direct PDF export.} The same plot at the same
  physical column width. The nominal font size is correct, though the
  glyph outlines can differ from Figure~\ref{fig:pgf}.}
  \label{fig:pdf}
\end{figure}

\section{Grouped scatter}
A categorical plot should not rely on hue alone.
Figure~\ref{fig:groups} assigns each category a distinct color and marker.
All observations in one category keep the same marker.

\begin{figure}[tb]
  \centering
  \includegraphics{generated/groups.pdf}
  \caption{Categorical scatter with redundant color and marker encoding,
  exported at the manuscript's actual column width.}
  \label{fig:groups}
\end{figure}

\section{Hybrid vector and raster export}
Large scatter clouds can make vector PDFs unnecessarily heavy.
Figure~\ref{fig:dense} rasterizes dense points but retains the surrounding
axes, labels and ticks as vectors.

\begin{figure}[tb]
  \centering
  \includegraphics{generated/dense.pdf}
  \caption{\textbf{Hybrid rendering.} Dense data points are rasterized,
  while simple graphic elements and typography remain vectorial.}
  \label{fig:dense}
\end{figure}

\section{Accessibility}
A color-vision simulation is a diagnostic approximation, not a clinical
prediction or universal certification. Figure~\ref{fig:access} compares
the same line plot in six views: original color, luminance-preserving
grayscale, print stress, deuteranopia, protanopia, and tritanopia.

\begin{figure*}[tp]
  \centering
  \includegraphics[width=\textwidth]{generated/accessibility.png}
  \caption{\textbf{Six-way comparison.} Original, black-and-white,
  degraded print, and three Machado-model color-vision-deficiency previews.
  The panel is rasterized; the source PDF and PGF above retain appropriate
  vector content.}
  \label{fig:access}
\end{figure*}

\begin{figure*}[tp]
  \centering
  \includegraphics[width=\textwidth]{generated/groups_accessibility.png}
  \caption{\textbf{Scatter accessibility panel.} The same categorical
  groups remain identified by their symbols in black and white and in
  simulated color-vision-deficiency conditions.}
  \label{fig:scatter-access}
\end{figure*}

\section{Reproducibility}
The accompanying script uses a fixed random seed, saves a JSON report
with font lookup, physical widths, and audit results, and compiles a
standalone proof confirming that the PDF figure's unscaled width
matches the measured LaTeX column.

Matching nominal point size, using PGF, and inspecting grayscale do not
eliminate a final visual review of dense legends, labels and scientific
clarity.

\end{document}
